\documentclass[11pt,a4paper]{article}

\usepackage[T1]{fontenc}
\usepackage[utf8]{inputenc}
\usepackage{lmodern}
\usepackage[margin=2.5cm]{geometry}
\usepackage{graphicx}
\usepackage{xcolor}
\usepackage{hyperref}
\usepackage{enumitem}
\usepackage{booktabs}
\usepackage{array}
\usepackage{microtype}

\graphicspath{{../Poster/images/}{Poster/images/}}

\hypersetup{
  colorlinks=true,
  linkcolor=blue,
  urlcolor=blue,
  citecolor=blue
}

\setlist[enumerate]{leftmargin=*, itemsep=0.35em, topsep=0.25em}
\setlist[itemize]{leftmargin=*, itemsep=0.25em, topsep=0.25em}
\setlength{\emergencystretch}{2em}

\begin{document}

\begin{center}
    {\LARGE \textbf{Quick Start Guide}\par}
    \vspace{0.2cm}
    \hrule
    \vspace{0.7cm}
    {\Large License Plate Detection System (Jetson Nano)\par}
    \vspace{0.6cm}
    \includegraphics[width=0.50\linewidth]{white.png}
\end{center}

\section{Introduction}
This quick start guide summarises the minimum steps needed to set up and run the License Plate Detection System in this repository.
The demo detects a German license plate in a live camera stream, runs OCR on the plate crop, and overlays the recognised text on the video output.

\section{Commissioning}

\subsection{Hardware Setup}
\textbf{Required hardware:} NVIDIA Jetson Nano (or compatible Linux PC), USB camera, monitor + HDMI cable, keyboard/mouse, and stable power supply.
For long-running operation we recommend an external SSD (USB 3.0) for logs/captures and (optional) swap.

\begin{enumerate}
    \item \textbf{Connect peripherals}\\
    Connect the USB camera, monitor (HDMI), keyboard/mouse, and (optional) external SSD to the Jetson Nano.

    \item \textbf{Power on and verify boot}\\
    Power on the Jetson Nano and log in to the desktop or terminal.

    \item \textbf{Place and aim the camera}\\
    Place the camera so the plate is readable (minimal angle, reduced glare). As a rule of thumb: the plate should occupy a noticeable portion of the image (avoid tiny plate crops).
\end{enumerate}

\textbf{The hardware is now ready for operation.}

\subsection{Software Setup (First Time Only)}
\textbf{Prerequisites:}
\begin{itemize}
    \item Python 3 and \texttt{pip}
    \item A working OpenCV camera backend (USB camera available as \texttt{/dev/video*} on Linux)
    \item Internet access at least once (recommended) to install Python packages and (optionally) OCR model files
\end{itemize}

\begin{enumerate}
    \item \textbf{Open a terminal and go to the project folder}\\
    From the repository root, change into the code folder:
    \begin{quote}\ttfamily
    cd Code/LicencePlateDetection
    \end{quote}

    \item \textbf{Create and activate a virtual environment}\\
    \begin{quote}\ttfamily
    python3 -m venv .venv\\
    source .venv/bin/activate
    \end{quote}

    \item \textbf{Install dependencies}\\
    \begin{quote}\ttfamily
    pip install --upgrade pip\\
    pip install -r requirements.txt
    \end{quote}

    \item \textbf{Ensure OCR assets are available (EasyOCR)}\\
    The OCR pipeline in \path{Code/LicencePlateDetection/jetson_app/ocr_pipeline.py} initialises EasyOCR with \texttt{download\_enabled=False}.
    If EasyOCR models are not already present on the device, OCR will fail to initialise.
    Recommended approach: run once on a machine with internet to populate \texttt{\textasciitilde/.EasyOCR/} and copy that folder to the Jetson Nano user account.
\end{enumerate}

\subsection{Starting the License Plate Detection Demo}
\begin{enumerate}
    \item \textbf{Activate the virtual environment}\\
    If not already active:
    \begin{quote}\ttfamily
    cd Code/LicencePlateDetection\\
    source .venv/bin/activate
    \end{quote}

    \item \textbf{Run the live demo}\\
    \begin{quote}\ttfamily
    python3 jetson\_app/main.py
    \end{quote}
    A window titled \textit{German Plate Detection + OCR} opens and shows the camera feed with detected plates and OCR text overlays.

    \item \textbf{Stop the demo}\\
    Press \texttt{q} in the video window to quit cleanly.
\end{enumerate}

\textbf{Output files:} The demo periodically saves annotated frames and plate crops to \path{Code/LicencePlateDetection/jetson_app/captures/}.

\subsection{Recommended Jetson Nano Configuration (Optional)}
For long-running gate/parking deployments, storage and memory pressure are common failure modes on Jetson-class devices.
If you use an external SSD, you can also place swap space on it to reduce out-of-memory crashes.

\begin{enumerate}
    \item \textbf{Mount an external SSD}\\
    Connect the SSD over USB 3.0 and mount it at a stable path (example: \texttt{/media/myssd}).
    Ensure the mount is persistent across reboots (e.g., via \texttt{/etc/fstab}).

    \item \textbf{(Optional) Create swap on the SSD}\\
    Example (creates a 6\,GB swap file):
    \begin{quote}\ttfamily
    sudo fallocate -l 6G /media/myssd/swapfile\\
    sudo chmod 600 /media/myssd/swapfile\\
    sudo mkswap /media/myssd/swapfile\\
    sudo swapon /media/myssd/swapfile
    \end{quote}
    Verify with \texttt{swapon --show}. To enable swap automatically at boot, add an entry to \texttt{/etc/fstab}.
\end{enumerate}

\section{Troubleshooting}
\begin{enumerate}
    \item \textbf{Camera not detected}\\
    If the program prints \texttt{ERROR: Camera not detected}, check the camera connection and that a video device exists.
    On Linux, try listing devices:
    \begin{quote}\ttfamily
    ls /dev/video*
    \end{quote}
    If your camera is not at index \texttt{0}, update \texttt{cv2.VideoCapture(0)} in:
    \begin{quote}\ttfamily
    Code/LicencePlateDetection/jetson\_app/detector.py
    \end{quote}

    \item \textbf{Model file not found}\\
    The detector loads the trained weights from:
    \begin{quote}\ttfamily
    Code/LicencePlateDetection/german\_lpr\_trained/train2/weights/best.pt
    \end{quote}
    If you moved files, restore the folder structure or update the model path in:
    \begin{quote}\ttfamily
    Code/LicencePlateDetection/jetson\_app/detector.py
    \end{quote}

    \item \textbf{OCR returns empty/unreadable text}\\
    Improve camera placement (focus, distance, angle), reduce glare, and ensure the plate crop is sharp.\\
    If OCR fails during startup, ensure EasyOCR model files exist under \texttt{\textasciitilde/.EasyOCR/} (see setup section).

    \item \textbf{Performance issues (low FPS / freezes)}\\
    Reduce camera resolution, reduce scene motion, and avoid running other heavy workloads.
    On Jetson-class devices, OCR is usually CPU-bound; minimise how often OCR runs if you customise the pipeline.
\end{enumerate}

\section{Optional: OCR Pipeline Test on Saved Plate Crops}
You can test OCR on a folder of plate crops (images) without running the live camera:
\begin{quote}\ttfamily
cd Code/LicencePlateDetection\\
source .venv/bin/activate\\
mkdir -p data/sample\_plates\\
python3 test\_ocr\_pipeline.py --dir data/sample\_plates
\end{quote}
Place plate-crop images (e.g., \texttt{.jpg}, \texttt{.png}) into \path{data/sample_plates} before running the test.

\end{document}
